\section{The rotation and the metric}\label{sec:rotation}

\subsection{Construction and exactness}

Complete $Q$ to an orthogonal matrix $R = [\,Q \mid Q_{\perp}\,] \in
\R^{N \times N}$. We use the Householder completion: $R$ is stored as $k$
reflectors, so applying $R$ or $R^{\top}$ to a vector costs $O(Nk)$ and no
$N \times N$ matrix is formed. Instead of $\Bstd$, sample
\begin{equation}\label{eq:rot}
\Bstd = R \begin{bmatrix} U \\ V \end{bmatrix} = Q U + Q_{\perp} V,
\qquad B = \Bstd L^{\top},
\end{equation}
where $Q \in \R^{N \times k}$, $Q_{\perp} \in \R^{N \times (N-k)}$, and
$U \in \R^{k \times q}$ and $V \in \R^{(N-k) \times q}$ have iid $\N(0,1)$
entries. $Q$, and hence $R$, is computed once from the design matrices
before sampling and is held fixed; it does not depend on any sampled
parameter.

\begin{proposition}[Exactness]\label{prop:exact}
For any fixed orthogonal $R$: (i) the map $(U, V) \mapsto \Bstd$ is linear
and volume-preserving; (ii) the prior it induces on $\Bstd$ is the original
iid $\N(0,1)$ prior; (iii) the posterior of $(B, \theta)$ is identical to
that of the non-centered parameterization (\ref{eq:ncp}), for every value of
$L$; and (iv) $U = Q^{\top} \Bstd$.
\end{proposition}

\begin{proof}
$\operatorname{vec}(\Bstd) = (I_q \otimes R)\operatorname{vec}([U; V])$
and $I_q \otimes R$ is orthogonal, which gives (i), and (ii) because
$(I_q \otimes R)\,\N(0, I_{Nq}) = \N(0, I_{Nq})$. The likelihood depends
on $(\Bstd, \theta)$ only, and their joint prior is unchanged, so the
posterior pushed forward to $(B, \theta)$ is unchanged: (iii). (iv)
follows from $Q^{\top} Q_{\perp} = 0$.
\end{proof}

Because the target is unchanged, NUTS and its warm-up adaptation are
unmodified; the rotation changes only which linear combinations of the
random effects are sampling coordinates, and therefore what an adapted
metric can represent.

Two choices in (\ref{eq:rot}) matter. First, the rotation acts on the
sampled, standardized coordinates: rotating centered effects $R^{\top} B$
would reintroduce the dependence on $\sigma_b$ that non-centering
removes. Second, $R$ acts on the subject index, the same for every column,
while $L$ acts on the column index. In vectorized form, $\operatorname{vec}(B) = (L \otimes I_N)
\operatorname{vec}(\Bstd)$ and $\operatorname{vec}(\Bstd) = (I_q \otimes
R)\operatorname{vec}([U; V])$, and
\[
(L \otimes I_N)(I_q \otimes R) = L \otimes R = (I_q \otimes R)(L \otimes I_N):
\]
the two factors act on different indices and commute, which is why the
construction is exact for every value of the covariance. A fixed rotation
aligning particular intercept-slope combinations with coordinate axes
would align them only at one value of $L$.

For a random intercept alone ($q = 1$), \texttt{jmjax} samples the
centered effects $b \sim \N(0, \sigma_b^2 I_N)$ directly, and the rotation
is applied to $b$. Because the covariance is a multiple of the identity,
$R^{\top} b$ has the same prior for every $\sigma_b$ and
Proposition~\ref{prop:exact} holds unchanged.

\subsection{Where the degeneracy goes}

\begin{lemma}[Localization]\label{lem:local}
Let $v \in \A_c$ with witness $d$. In the coordinates
$(\beta, U, V)$, the likelihood-invariant shift $(\beta, B e_c) \mapsto
(\beta + d, B e_c + v)$ is
\[
\beta \mapsto \beta + d, \qquad
U \mapsto U + (Q^{\top} v)\, e_c^{\top} (L^{\top})^{-1}, \qquad
V \mapsto V,
\]
for every value of $L$. Every such shift therefore lies in the coordinate
block $(\beta, U)$, of dimension $p + kq$.
\end{lemma}

\begin{proof}
Adding $v$ to column $c$ of $B = \Bstd L^{\top}$ adds $\Delta = v\,
e_c^{\top} (L^{\top})^{-1}$ to $\Bstd$. Every column of $\Delta$ is a
multiple of $v \in \spn(Q)$, so $Q_{\perp}^{\top} \Delta = 0$ and
$Q^{\top} \Delta = (Q^{\top} v)\, e_c^{\top} (L^{\top})^{-1}$.
\end{proof}

Without the rotation, the same shift has a component on every one of the
$N$ rows of $\Bstd$. With it, the invariant directions mix the columns of
$U$ in proportions set by the sampled $L$, but stay inside the fixed block
$(\beta, U)$; rotating the intercept column alone would capture them only
for uncorrelated random effects. $U$ has $kq$ coordinates, a few of which
may not be invariant for their own column (in the design of
Section~\ref{sec:simulation}, 4 of the 6 are invariant and 2 are not); including them in the dense
block costs little.

\subsection{The ridge, and the dense block}

Along each invariant direction the priors still locate the split, so the
posterior is proper but weakly identified. The next result describes the
ridge in one plane.

\begin{proposition}[The ridge]\label{prop:ridge}
Consider one invariant direction in the reduced model that holds every
other coordinate fixed. Let $\tilde\beta$ be the coordinate of $\beta$ along
the witness, with a flat prior, and $u$ the matching coordinate of $U$,
with prior $\N(0,1)$, and suppose the likelihood informs only $s =
\tilde\beta + \kappa u$, with precision $\tau$. Then
\begin{enumerate}
\item[(i)] $\corr(\tilde\beta, u \mid y) = -\{1 + 1/(\tau\kappa^2)\}^{-1/2}$;
\item[(ii)] under the ideal diagonal metric, equal to the exact marginal
posterior variances (the quantity warm-up adaptation estimates), the
whitened covariance of the plane has eigenvalues $1 \pm |r|$, $r$ the
correlation in (i). Its wide direction has variance at most 2, but its
narrow direction has variance $1 - |r| \approx 1/(2\tau\kappa^2)$;
\item[(iii)] a dense metric over $(\tilde\beta, u)$ whitens the plane
exactly;
\item[(iv)] for the mean direction of the random intercepts, $\kappa =
\sigma_0/\sqrt{N}$ and $\tau = \sum_i n_i/\sigma_e^2$, so $\tau\kappa^2 =
\bar n\,\sigma_0^2/\sigma_e^2$, the information each subject carries about
its own level. It does not depend on $N$.
\end{enumerate}
\end{proposition}

\begin{proof}
The posterior precision of $(\tilde\beta, u)$ is $\tau (1, \kappa)^{\top}
(1, \kappa) + \operatorname{diag}(0, 1)$, whose inverse is
\[
\tau^{-1}\begin{pmatrix} 1 + \tau\kappa^2 & -\tau\kappa \\ -\tau\kappa
& \tau \end{pmatrix};
\]
this gives (i). Whitening a $2 \times 2$ covariance
by its diagonal gives its correlation matrix, whose eigenvalues are $1 \pm
|r|$: (ii). (iii) is the definition of a dense metric. For (iv), moving
the mean direction by $u$ moves every subject's level by $\sigma_0
u/\sqrt{N}$, and the data inform a common shift of all levels with
precision $\sum_i n_i / \sigma_e^2$.
\end{proof}

Heuristically, for a Gaussian target the narrow direction bounds the
stable leapfrog step size and the wide one sets the trajectory length, so
the number of leapfrog steps per independent draw grows like
$\{(1 + |r|)/(1 - |r|)\}^{1/2} \approx 2\sqrt{\tau\kappa^2}$
\citep{neal2011}. Proposition~\ref{prop:ridge} is exact for the mean
intercept direction, whose witness is $\beta_0$; for other directions the
witness mixes several coordinates of $\beta$, and (ii)--(iii) describe the
whole $(\beta, U)$ block. The rotation alone thus concentrates the
direction in a few coordinates, but under a diagonal metric it is still
narrow, most so with many visits per subject or little noise. A dense
block over $(\beta, U)$ whitens it. Without the rotation, the same fix
needs a dense metric over all $Nq$ random-effect coordinates, with
$O(N^2q^2)$ storage and work per metric-vector product;
Section~\ref{sec:generality} measures this alternative.

Figure~\ref{fig:ridge} shows this in the exact Gaussian posterior of the
longitudinal submodel with the variance components fixed at the
simulation values of Section~\ref{sec:simulation}. The conditional
correlation of $\beta_0$ and the mean intercept direction is $-0.97$
unrotated and still $-0.97$ after rotation with a diagonal metric; the
dense block brings it to $-0.01$. The condition number of the whole
whitened posterior, the ratio of the extreme eigenvalues of $M^{-1/2}
\Sigma M^{-1/2}$ with $\Sigma$ the exact posterior covariance and $M$ the
metric ($\operatorname{diag}\Sigma$, or $\operatorname{diag}\Sigma$ with
its $(\beta, U)$ block replaced by that of $\Sigma$), falls from 4,521 to
368, and by factors of 5 to 70 across the designs examined; what remains
is conditioning any fit has, such as the correlation between each
subject's own intercept and slope.

\begin{figure}[t]
\centering
\includegraphics[width=\textwidth]{fig1_ridge.pdf}
\caption{The location ridge and what removes it. Exact conditional
posterior of $\beta_0$ and the mean intercept direction $u$ in the units
the HMC metric sees them (Gaussian longitudinal submodel, variance
components fixed; $N = 300$, 2--11 visits, $\rho = 0.3$, $\sigma_e =
0.3$). Solid: 1, 2 and 3 standard deviation contours; dashed: the 2
standard deviation contour of a perfectly scaled target. (a) Unrotated,
diagonal metric. (b) Rotated, diagonal metric: $u$ is a single coordinate
but the ridge remains. (c) Rotated, dense block over $(\beta, U)$.
``cond.'': condition number of the whole whitened posterior.}
\label{fig:ridge}
\end{figure}

\subsection{What a fixed metric cannot do}\label{sec:residual}

By Lemma~\ref{lem:local}, the ridge's orientation depends on the sampled
$L$: on $\sigma_0$ through $\kappa \propto \sigma_0$ and, for directions
shared between the columns, on $\rho$, while NUTS adapts one fixed metric.
For the $\sigma_0$ part the loss can be estimated. Let $\delta$ be the
posterior coefficient of variation of $\sigma_0$, roughly $1/\sqrt{2N}$ in
an idealized variance-component model with $N$ groups (we use that value
only to set the scale). A relative change
$\delta$ in $\kappa$ tilts the plane's narrow direction, in the units of a
metric fitted at the original value, by an angle of order
$\delta\sqrt{\tau\kappa^2}$, so a metric that whitens the plane exactly at
one value of $\sigma_0$ leaves a condition number of $1 +
O(\delta\sqrt{\tau\kappa^2})$ elsewhere. With $\delta
\approx (2N)^{-1/2}$ and $\tau\kappa^2 = \bar n\,\sigma_0^2/\sigma_e^2$
this is
\[
1 + O\Bigl(\sqrt{\bar n\,\sigma_0^2 / (\sigma_e^2 N)}\Bigr),
\]
which is small when there are many subjects relative to the information
each carries, the usual regime for longitudinal data. In the exact
Gaussian model of Figure~\ref{fig:ridge} the worst block condition number
over $\sigma_0 (1 \pm \delta)$ was 1.76, 1.50 and 1.31 for $N = 150$, 300
and 600 with 2--11 visits, and 2.71, 2.05 and 1.64 with 8--20 visits. The
analysis does not cover the $\rho$ part, which grows as $|\rho| \to 1$;
Section~\ref{sec:stress} stress-tests both regimes. For $q = 1$ in
centered coordinates only the scale of the rotated coordinate's prior
involves $\sigma_b$, a mismatch of order $\delta$.

\subsection{Implementation and cost}

\texttt{jmjax} implements (\ref{eq:rot}) in NumPyro: $U$ and $V$ are the
sampled sites, $\Bstd$ is computed from them with $q$ applications of the
$k$ reflectors, and the mass matrix is dense on $(\beta, U)$ and diagonal
elsewhere, using NumPyro's dense blocks over groups of sites, adapted
during warm-up as usual. The extra cost per gradient is $O(Nkq)$. In the
design of Section~\ref{sec:simulation} the dense block has $p + kq = 10$
coordinates, and wall time per fit was 0.83--1.10 times that of the
unrotated fit. The rotation and the dense block are the package's default
for NUTS fits with a random intercept only, or with $q \geq 2$ correlated
random effects.
